11.0 ODPIERANIE CZARÓW

Ogólna zasada:

Jeżeli na jakąś postać lub potwora rzucany jest czar, użyty jest przeciwko niej lub jemu magiczny przedmiot lub jeżeli postać bada pewne znaleziska, to może zaistnieć potrzeba podjęcia próby odparcia czaru, a tym samym zniwelowania jego efektów.

Sposób postępowania:

Każda postać oraz każdy potwór posiada współczynnik odporność.
Gdy istnieje konieczność lub możliwość odparcia czaru gracz rzuca 1K6 i porównuje liczbę wyrzuconych oczek z wielkością tego współczynnika u postaci, która ma czar odeprzeć.
Jeżeli wyrzucona liczba oczek jest równa lub mniejsza od wielkości współczynnika odporność czar został odparty.

ustalić ile razy ten czar mógł być użyty podczas bieżącej przygody.
Po wyczerpaniu limitu, określonego przez liczbę wyrzuconych oczek postać nie może już rzucać czaru uzdrowienia (jest to wyjątek od zasady 10.2).

Odmłodzenia (koszt: 2) – Ten czar ma taki sam efekt jak uzdrowienie, z tą różnicą, iż leczy z 1K6+1 ran.
Również może być użyty tylko ograniczoną liczbę razy w czasie jednej przygody.
(Ograniczenie według zasady jak przy uzdrawieniu).

Delikatne Palce (koszt: 1) – Ten czar zwiększa zdolności ,,pułapki'' o 3 punkty, ale tylko na czas potrzebny na przeprowadzenie jednej próby unieszkodliwienia pułapki.
Nie może być rzucony na postać, która w ogóle nie posiada zdolności ,,pułapki''.

[10.7] Czary, pomagające w negocjacjach (typ N) mogą być używane przed rzutem kostką, przesądzającym o wyniku negocjacji, zaś czary typu V – podczas próby wykupu.

Elokwencja (koszt: 1) – Ten czar pozwala na dodanie 4 punktów do liczby oczek, wyrzuconej podczas negocjacji (dodatkowo, niezależnie od wszelkich modyfikacji, wynikających ze zdolności negocjujące postaci).

Zastraszenie (koszt: 2) – Daje efekt (bez konieczności rzucania kostką) analogiczny do ,,zastraszenia'' w negocjacjach (patrz 8.6).

Zniechęcenie (koszt: 1) – Ten czar daje podobny efekt jak zastraszenie, z tym, że potwór oddaje wszystkie swoje skarby.

Sugestia (koszt: 1) – Pozwala oddziałowi na odjęcie 2 punktów od liczby oczek, wyrzuconych przy próbie wykupienia się.

Pochlebstwo (koszt: 2) – Pozwala oddziałowi na dodanie 4 punktów do liczby oczek wyrzuconych przy próbie wykupienia się.

[10.8] Umiejętność rzucania czarów specjalnych (typ S) może być zdobyta tylko w toku gry, tzn. nie może one być wybierane na początku, do ,,wstępnej puli'' czarów.

Gniew Boży (koszt: NIECZYTELNE) – Ten czar może być użyty tylko podczas walki.
Umiejętność jego rzucania jest nabywana poprzez znalezienie manuskryptu lub ołtarza Malthusa (patrz 13.9).
Jeżeli nie zostanie odparty zadaje 2K6+2 ran.

Wskrzeszenie (koszt: 5) – Tego czaru nie wolno używać podczas walki.
Umiejętność jego rzucania może być nabyta tylko przez znalezienie pierścienia ,,Wskrzeszenie'' (patrz 14.9) lub regału z książkami (patrz 13.9).
Może być użyty do wskrzeszenia postaci, która zginęła podczas walki.
Musi być rzucony natychmiast po zakończeniu walki, w której postać ta poniosła śmierć, gdyż użyty później nie będzie działał.
Wskrzeszona postać ma tyle ran, ile miała w momencie wejścia do labiryntu.

[10.9] Tabela ,,Zestawienie czarów'' – patrz plansza z tabelami.